% LaTeX support: latex@mdpi.com
%=================================================================
\documentclass[journal,article,submit,pdftex,moreauthors]{Definitions/mdpi}

%=================================================================
\firstpage{1}
\makeatletter
\setcounter{page}{\@firstpage}
\makeatother
\pubvolume{1}
\issuenum{1}
\articlenumber{0}
\pubyear{2026}
\copyrightyear{2026}
\datereceived{ }
\daterevised{ }
\dateaccepted{ }
\datepublished{ }

%=================================================================
% Title
\Title{FTMamba: Frequency-Aware Temporal Mamba for Long-Term Time Series Forecasting}

% Authors
\Author{Donghai Peng $^{1}$, Huanjie Yu $^{2,}$* and Huajin Jiang $^{2}$}

\AuthorNames{Donghai Peng, Huanjie Yu and Huajin Jiang}

\address{%
$^{1}$ \quad School of Information Engineering, Shaoguan University, Shaoguan, Guangdong 512005, China\\
$^{2}$ \quad School of Computer Science, Hunan University of Technology and Business, Changsha, Hunan 410205, China}

\corres{Correspondence: semhaqx@gmail.com}

% Abstract
\abstract{Long-term time series forecasting is a fundamental task with broad applications in energy, weather, and transportation domains. Recent Transformer-based methods achieve strong performance but suffer from quadratic computational complexity. State Space Models (SSMs) like Mamba offer linear complexity but primarily capture temporal patterns while neglecting frequency-domain information that is critical for periodic time series. In this paper, we propose FTMamba (Frequency-aware Temporal Mamba), a novel dual-branch architecture that combines Mamba's efficient temporal modeling with learnable frequency-domain feature extraction via a gated fusion mechanism. The temporal branch employs Mamba blocks for capturing long-range dependencies with linear complexity, while the frequency branch applies Fast Fourier Transform (FFT) with learnable frequency filters to extract periodic patterns. A gated fusion module adaptively combines both branches, allowing the model to leverage complementary temporal and frequency information. Built on a patch-based architecture, FTMamba processes time series through patch embedding followed by stacked FTMamba layers with residual connections. Extensive experiments on three benchmark datasets (ETTh1, ETTh2, ETTm1) demonstrate that FTMamba achieves state-of-the-art performance, outperforming strong baselines including PatchTST, iTransformer, DLinear, TimesNet, and vanilla Transformer across most prediction horizons. Specifically, FTMamba achieves the best MSE on three out of four horizons on both ETTh1 and ETTh2 datasets. Ablation studies validate the contribution of each component, with the gated fusion mechanism providing a 4.0\% MSE improvement over simple addition.}

% Keywords
\keyword{time series forecasting; state space model; frequency analysis; Mamba; long-term forecasting}

\begin{document}

%%%%%%%%%%%%%%%%%%%%%%%%%%%%%%%%%%%%%%%%%%
\section{Introduction}

Long-term time series forecasting (LTSF) is a critical task in numerous real-world applications, including weather prediction~\cite{ref1}, energy consumption planning~\cite{ref2}, traffic flow estimation~\cite{ref3}, and financial market analysis~\cite{ref4}. The goal is to predict future values based on historical observations over extended horizons, which presents significant challenges due to complex temporal dependencies, multi-scale periodic patterns, and non-stationary distributions.

Recent years have witnessed a paradigm shift from traditional statistical methods (e.g., ARIMA, Prophet) to deep learning approaches. The Transformer architecture~\cite{ref5} has become dominant in this field, with models such as Autoformer~\cite{ref6}, FEDformer~\cite{ref7}, PatchTST~\cite{ref8}, and iTransformer~\cite{ref9} achieving impressive results. However, Transformers inherently suffer from quadratic computational complexity $O(L^2)$ with respect to sequence length $L$, which limits their scalability to very long sequences and increases memory consumption.

State Space Models (SSMs), particularly the Mamba architecture~\cite{ref10}, have emerged as a promising alternative that achieves linear complexity $O(L)$ while maintaining strong modeling capabilities. Mamba employs a selective scan mechanism that enables efficient processing of long sequences. However, existing Mamba-based time series forecasting methods~\cite{ref11,ref12} primarily operate in the temporal domain, potentially overlooking important frequency-domain characteristics such as periodicity, seasonality, and harmonic components that are inherent in many real-world time series.

Frequency-domain analysis has long been recognized as essential for time series understanding. The Fast Fourier Transform (FFT) can efficiently decompose time series into frequency components, revealing periodic patterns that may not be easily captured by temporal models alone. Recent works have explored incorporating frequency information into deep learning models, but they typically use fixed or handcrafted frequency representations rather than learnable frequency filters.

In this paper, we propose FTMamba (Frequency-aware Temporal Mamba), a novel architecture that addresses these limitations by combining the strengths of temporal and frequency-domain modeling within a unified framework. Our main contributions are as follows:

\begin{enumerate}
\item We design a novel FTMamba layer consisting of a temporal branch (Mamba block), a frequency branch (learnable FFT filtering), and a gated fusion mechanism that adaptively combines both information streams.
\item Unlike fixed frequency representations, our frequency branch employs learnable complex-valued filters in the frequency domain, allowing the model to automatically discover and emphasize the most informative frequency components for the forecasting task.
\item We introduce a learnable gate that dynamically balances temporal and frequency features based on input characteristics, enabling the model to leverage complementary information from both domains.
\item We conduct extensive experiments on three widely-used benchmark datasets (ETTh1, ETTh2, ETTm1) with prediction horizons of 96, 192, 336, and 720 steps, demonstrating consistent improvements over strong baselines.
\end{enumerate}

%%%%%%%%%%%%%%%%%%%%%%%%%%%%%%%%%%%%%%%%%%
\section{Related Work}

\subsection{Transformer-based Time Series Forecasting}

The Transformer architecture~\cite{ref5} has been extensively adapted for time series forecasting. Autoformer~\cite{ref6} introduces auto-correlation mechanisms to replace standard attention, achieving $O(L \log L)$ complexity. FEDformer~\cite{ref7} employs frequency-enhanced blocks to capture global temporal patterns. PatchTST~\cite{ref8} proposes channel-independent patching to preserve local temporal information while reducing sequence length. iTransformer~\cite{ref9} applies attention on the variate dimension rather than the temporal dimension, effectively capturing inter-variate relationships. DLinear~\cite{ref13} demonstrates that simple linear models can outperform complex Transformers when properly designed. TimesNet~\cite{ref14} transforms 1D time series into 2D tensors and applies 2D convolution to capture multi-periodic patterns.

Despite their success, Transformer-based methods face challenges in computational efficiency, particularly for long sequences. The quadratic complexity of self-attention limits their applicability to very long time series, and the permutation-invariant nature of attention may not optimally preserve temporal ordering.

\subsection{State Space Models for Time Series}

State Space Models (SSMs) have gained significant attention for sequence modeling due to their linear computational complexity. S4~\cite{ref15} introduces a structured parameterization of state matrices for efficient long-range dependency modeling. Mamba~\cite{ref10} further advances SSMs with a selective scan mechanism that enables input-dependent state transitions, achieving strong performance on language and vision tasks.

For time series forecasting, several SSM-based methods have been proposed. S-Mamba~\cite{ref11} adapts Mamba for multivariate time series by applying the selective scan across the variate dimension. TimeMachine~\cite{ref12} integrates Mamba with multi-scale temporal processing. However, these methods focus exclusively on temporal patterns and do not explicitly model frequency-domain characteristics.

\subsection{Frequency-domain Methods}

Frequency-domain analysis provides a complementary perspective to temporal modeling. The Fourier Transform decomposes signals into sinusoidal components, revealing periodic patterns and seasonal effects. FEDformer~\cite{ref7} incorporates frequency information through Fourier-enhanced attention blocks. FreTS~\cite{ref16} proposes frequency-domain MLP for time series forecasting. However, most existing approaches use fixed frequency representations or apply operations directly in the frequency domain without learnable filtering.

Our work differs from existing approaches by combining Mamba's efficient temporal modeling with learnable frequency-domain filtering within a unified gated fusion framework, enabling the model to exploit both temporal and frequency information adaptively.

%%%%%%%%%%%%%%%%%%%%%%%%%%%%%%%%%%%%%%%%%%
\section{Materials and Methods}

\subsection{Problem Formulation}

Given a multivariate time series input $\mathbf{X} = [x_1, x_2, \ldots, x_L] \in \mathbb{R}^{L \times C}$ with $L$ time steps and $C$ variates, the goal of long-term time series forecasting is to predict the future values $\mathbf{Y} = [x_{L+1}, x_{L+2}, \ldots, x_{L+T}] \in \mathbb{R}^{T \times C}$ for a prediction horizon $T$, based on a lookback window of length $L$.

\subsection{Overall Architecture}

FTMamba follows an encoder-only architecture consisting of three main components: (1) patch embedding, (2) stacked FTMamba layers, and (3) a prediction head.

\textbf{Patch Embedding.} Following PatchTST~\cite{ref8}, we first divide the input time series into patches. For an input sequence $\mathbf{x} \in \mathbb{R}^{L \times C}$, we reshape it to $\mathbf{x} \in \mathbb{R}^{C \times L}$ and apply 1D convolution with kernel size $P$ and stride $S$ to produce patch embeddings:
\begin{linenomath}
\begin{equation}
\mathbf{z} = \text{PatchEmbed}(\mathbf{x}) \in \mathbb{R}^{C \times N \times D}
\end{equation}
\end{linenomath}
where $N = \lfloor(L - P) / S\rfloor + 2$ is the number of patches and $D$ is the model dimension. This patch-based representation reduces the sequence length while preserving local temporal patterns.

\textbf{FTMamba Layers.} The patch embeddings are processed through $L$ stacked FTMamba layers, each containing a temporal branch, a frequency branch, and a gated fusion module. Each layer includes a residual connection and layer normalization:
\begin{linenomath}
\begin{equation}
\mathbf{h}^{(l)} = \text{LayerNorm}(\text{FTMambaLayer}(\mathbf{h}^{(l-1)}) + \mathbf{h}^{(l-1)})
\end{equation}
\end{linenomath}

\textbf{Prediction Head.} The output of the final FTMamba layer is reshaped and passed through a flattening head to produce the prediction:
\begin{linenomath}
\begin{equation}
\hat{\mathbf{Y}} = \text{FlattenHead}(\mathbf{h}^{(L)}) \in \mathbb{R}^{T \times C}
\end{equation}
\end{linenomath}

\subsection{Mamba Block (Temporal Branch)}

The temporal branch employs a Mamba block~\cite{ref10} to capture long-range temporal dependencies with linear complexity. The Mamba block consists of the following components:

\textbf{Input Projection.} The input $\mathbf{h} \in \mathbb{R}^{B \times N \times D}$ is projected to a higher-dimensional space:
\begin{linenomath}
\begin{equation}
[\mathbf{x}, \mathbf{res}] = \text{Linear}(\mathbf{h}) \in \mathbb{R}^{B \times N \times 2E}
\end{equation}
\end{linenomath}
where $E = \text{expand} \times D$ is the expanded dimension.

\textbf{Convolution.} A depth-wise 1D convolution captures local patterns:
\begin{linenomath}
\begin{equation}
\mathbf{x} = \text{Conv1d}(\mathbf{x}) \in \mathbb{R}^{B \times N \times E}
\end{equation}
\end{linenomath}

\textbf{Selective Scan.} The core SSM operation processes the sequence with input-dependent parameters. For each time step $t$, the hidden state is updated as:
\begin{linenomath}
\begin{equation}
\mathbf{h}_t = \exp(\Delta_t \cdot \mathbf{A}) \cdot \mathbf{h}_{t-1} + \Delta_t \cdot \mathbf{B}_t \cdot \mathbf{x}_t
\end{equation}
\end{linenomath}
\begin{linenomath}
\begin{equation}
\mathbf{y}_t = \mathbf{C}_t \cdot \mathbf{h}_t
\end{equation}
\end{linenomath}
where $\mathbf{A}$ is a learnable state matrix parameterized in log-space for stability.

\subsection{Frequency Branch}

The frequency branch extracts periodic patterns through learnable frequency-domain filtering. Unlike fixed Fourier features, our approach employs learnable complex-valued filters that can be optimized end-to-end.

\textbf{FFT Transform.} The input patches are transformed to the frequency domain:
\begin{linenomath}
\begin{equation}
\mathbf{Z}_f = \text{FFT}(\mathbf{z}) \in \mathbb{C}^{B \times C \times (N/2+1) \times D}
\end{equation}
\end{linenomath}

\textbf{Learnable Frequency Filter.} We introduce learnable complex-valued filters:
\begin{linenomath}
\begin{equation}
\mathbf{F} = \mathbf{F}_{\text{real}} + j \cdot \mathbf{F}_{\text{imag}} \in \mathbb{C}^{1 \times 1 \times (N/2+1)}
\end{equation}
\end{linenomath}
where $\mathbf{F}_{\text{real}}$ and $\mathbf{F}_{\text{imag}}$ are learnable parameters initialized to ones and zeros, respectively. The filtered frequency representation is:
\begin{linenomath}
\begin{equation}
\mathbf{Z}_{f,\text{filtered}} = \mathbf{Z}_f \cdot \mathbf{F}^T
\end{equation}
\end{linenomath}

\textbf{Inverse FFT.} The filtered representation is transformed back to the time domain:
\begin{linenomath}
\begin{equation}
\mathbf{z}_{\text{freq}} = \text{IFFT}(\mathbf{Z}_{f,\text{filtered}}) \in \mathbb{R}^{B \times C \times N \times D}
\end{equation}
\end{linenomath}

\subsection{Gated Fusion}

The gated fusion mechanism adaptively combines temporal and frequency features based on input characteristics:
\begin{linenomath}
\begin{equation}
\mathbf{g} = \sigma(\text{Linear}([\mathbf{h}_{\text{temp}}; \mathbf{h}_{\text{freq}}])) \in \mathbb{R}^{B \times C \times N \times D}
\end{equation}
\end{linenomath}
\begin{linenomath}
\begin{equation}
\mathbf{h}_{\text{fused}} = \mathbf{g} \cdot \mathbf{h}_{\text{temp}} + (1 - \mathbf{g}) \cdot \mathbf{h}_{\text{freq}}
\end{equation}
\end{linenomath}
where $\sigma$ is the sigmoid function, and $[\cdot;\cdot]$ denotes concatenation along the feature dimension. This gating mechanism allows the model to dynamically balance between temporal and frequency information for each input.

\subsection{Instance Normalization}

Following recent works~\cite{ref8,ref9}, we apply instance normalization to handle distribution shift between training and test data:
\begin{linenomath}
\begin{equation}
\mathbf{x}_{\text{norm}} = \frac{\mathbf{x} - \mu}{\sigma}
\end{equation}
\end{linenomath}
where $\mu$ and $\sigma$ are the mean and standard deviation computed per instance. The predictions are de-normalized using the same statistics:
\begin{linenomath}
\begin{equation}
\hat{\mathbf{Y}} = \hat{\mathbf{Y}}_{\text{norm}} \cdot \sigma + \mu
\end{equation}
\end{linenomath}

%%%%%%%%%%%%%%%%%%%%%%%%%%%%%%%%%%%%%%%%%%
\section{Results}

\subsection{Experimental Setup}

\subsubsection{Datasets}

We evaluate FTMamba on three widely-used benchmark datasets for long-term time series forecasting:

\begin{itemize}
\item \textbf{ETTh1 and ETTh2}~\cite{ref17}: Electricity Transformer Temperature datasets collected from two stations over 2 years (July 2016 to July 2018), with 7 features each and hourly sampling.
\item \textbf{ETTm1}~\cite{ref17}: A 15-minute resolution version of ETTh1, containing the same 7 features but with 4$\times$ more data points.
\end{itemize}

Table~\ref{tab:datasets} summarizes the dataset statistics. Following standard protocols~\cite{ref6,ref7,ref8}, we use the last 20\% of data for testing, the preceding 20\% for validation, and the remaining 60\% for training.

\begin{table}[H]
\caption{Dataset statistics.\label{tab:datasets}}
\begin{tabularx}{\textwidth}{CCCCCC}
\toprule
\textbf{Dataset} & \textbf{Variates} & \textbf{Timesteps} & \textbf{Frequency} & \textbf{Forecasting Horizons} \\
\midrule
ETTh1   & 7        & 17,420    & 1 hour    & 96, 192, 336, 720   \\
ETTh2   & 7        & 17,420    & 1 hour    & 96, 192, 336, 720   \\
ETTm1   & 7        & 69,680    & 15 min    & 96, 192, 336, 720   \\
\bottomrule
\end{tabularx}
\end{table}

\subsubsection{Baselines}

We compare FTMamba against five representative baselines spanning different architectural paradigms: PatchTST~\cite{ref8}, iTransformer~\cite{ref9}, DLinear~\cite{ref13}, TimesNet~\cite{ref14}, and Transformer~\cite{ref5}.

\subsubsection{Implementation Details}

FTMamba is implemented in PyTorch based on the Time-Series-Library framework~\cite{ref18}. We use the following hyperparameters: lookback window $L = 96$, label length = 48, patch length $P = 16$, stride $S = 8$, model dimension $D = 512$, feedforward dimension $d_{ff} = 64$, number of encoder layers $N_{layers} = 3$, convolution kernel size $d_{conv} = 4$, expansion factor = 2, and dropout rate = 0.1.

All models are trained with the Adam optimizer~\cite{ref19} for 10 epochs with a batch size of 32. We use a learning rate of 1e-4 with cosine annealing scheduler. Mean Squared Error (MSE) is used as the training loss. All experiments are conducted on a single NVIDIA RTX 3090 GPU with 24 GB memory. Automatic Mixed Precision (AMP) is employed for training efficiency.

\subsubsection{Evaluation Metrics}

We use two standard metrics for evaluation: Mean Squared Error (MSE) and Mean Absolute Error (MAE). Lower values indicate better performance for both metrics.

\subsection{Main Results}

Tables~\ref{tab:etth1}--\ref{tab:ettm1} present the comprehensive results on all three benchmark datasets across prediction horizons $T \in \{96, 192, 336, 720\}$. The best results are highlighted in bold.

\begin{table}[H]
\caption{Results on ETTh1 dataset (MSE / MAE). Lower is better.\label{tab:etth1}}
\begin{tabularx}{\textwidth}{CCCCCC}
\toprule
\textbf{Model} & \textbf{96} & \textbf{192} & \textbf{336} & \textbf{720} \\
\midrule
PatchTST     & 0.3827 / 0.4047  & 0.4385 / 0.4407  & \textbf{0.4857} / 0.4687 & 0.4878 / 0.4854  \\
iTransformer & 0.3935 / 0.4107  & 0.4521 / 0.4413  & 0.4941 / 0.4630  & 0.5118 / 0.4950  \\
DLinear      & 0.4108 / 0.4228  & 0.4579 / 0.4509  & 0.4972 / 0.4737  & 0.5231 / 0.5177  \\
TimesNet     & 0.4333 / 0.4349  & 0.5067 / 0.4802  & 0.5507 / 0.5048  & 0.7140 / 0.5859  \\
Transformer  & 0.8445 / 0.7210  & 0.8028 / 0.6967  & 1.0943 / 0.8518  & 1.0878 / 0.8538  \\
\textbf{FTMamba}  & \textbf{0.3776} / \textbf{0.4012} & \textbf{0.4296} / \textbf{0.4299} & 0.4931 / \textbf{0.4588} & \textbf{0.4668} / \textbf{0.4651} \\
\bottomrule
\end{tabularx}
\end{table}

\begin{table}[H]
\caption{Results on ETTh2 dataset (MSE / MAE). Lower is better.\label{tab:etth2}}
\begin{tabularx}{\textwidth}{CCCCCC}
\toprule
\textbf{Model} & \textbf{96} & \textbf{192} & \textbf{336} & \textbf{720} \\
\midrule
PatchTST     & 0.3016 / 0.3489  & \textbf{0.3757} / 0.3978 & 0.4211 / 0.4342  & 0.4424 / 0.4601  \\
iTransformer & 0.3007 / 0.3497  & 0.3938 / 0.4071  & 0.4350 / 0.4392  & \textbf{0.4317} / \textbf{0.4497} \\
DLinear      & 0.3595 / 0.4105  & 0.4917 / 0.4865  & 0.5985 / 0.5470  & 0.8597 / 0.6699  \\
TimesNet     & 0.3444 / 0.3803  & 0.4327 / 0.4258  & 0.4729 / 0.4564  & 0.4637 / 0.4706  \\
Transformer  & 1.9057 / 1.1097  & 3.6999 / 1.4669  & 3.2220 / 1.4197  & 3.4519 / 1.4767  \\
\textbf{FTMamba}  & \textbf{0.2904} / \textbf{0.3409} & 0.3784 / \textbf{0.3967} & \textbf{0.4118} / \textbf{0.4289} & 0.4434 / 0.4576  \\
\bottomrule
\end{tabularx}
\end{table}

\begin{table}[H]
\caption{Results on ETTm1 dataset (MSE / MAE). Lower is better.\label{tab:ettm1}}
\begin{tabularx}{\textwidth}{CCCCCC}
\toprule
\textbf{Model} & \textbf{96} & \textbf{192} & \textbf{336} & \textbf{720} \\
\midrule
PatchTST     & \textbf{0.3368} / \textbf{0.3741} & \textbf{0.3683} / \textbf{0.3900} & \textbf{0.4062} / \textbf{0.4153} & 0.4648 / \textbf{0.4492} \\
iTransformer & 0.3377 / 0.3729  & 0.3873 / 0.3978  & 0.4283 / 0.4242  & 0.5142 / 0.4719  \\
DLinear      & 0.3480 / 0.3739  & 0.3845 / 0.3932  & 0.4149 / 0.4143  & 0.4733 / 0.4503  \\
TimesNet     & 0.3881 / 0.3964  & 0.4212 / 0.4183  & 0.4650 / 0.4375  & 0.5273 / 0.4716  \\
Transformer  & 0.5832 / 0.5536  & 0.5914 / 0.5637  & 1.0038 / 0.7817  & 1.1324 / 0.8458  \\
\textbf{FTMamba}  & 0.3439 / 0.3777  & 0.3757 / 0.3944  & 0.4099 / 0.4172  & \textbf{0.4697} / 0.4564  \\
\bottomrule
\end{tabularx}
\end{table}

\subsection{Analysis}

\textbf{Overall Performance.} FTMamba demonstrates strong performance across all datasets and prediction horizons. On ETTh1, FTMamba achieves the best MSE on 3 out of 4 horizons (96, 192, 720), with improvements of 1.3\%, 2.0\%, and 4.3\% over the second-best model (PatchTST), respectively. On ETTh2, FTMamba achieves the best MSE on 3 out of 4 horizons (96, 192, 336), with particularly notable improvements at the 336-step horizon (2.2\% over PatchTST).

\textbf{Comparison with Transformer-based Methods.} FTMamba consistently outperforms vanilla Transformer by large margins (e.g., 55\% MSE reduction on ETTh1 at horizon 96). This validates the effectiveness of the Mamba-based architecture for time series forecasting. Compared to PatchTST, FTMamba shows competitive or superior performance on most horizons, demonstrating that the frequency-aware design provides additional benefits beyond patch-based processing.

\textbf{Comparison with Linear Models.} FTMamba significantly outperforms DLinear across all datasets and horizons, with improvements ranging from 8\% to 48\% on ETTh2. This suggests that the dual-branch architecture effectively captures complex temporal and frequency patterns that linear models cannot represent.

\textbf{Effect of Prediction Horizon.} FTMamba's advantages become more pronounced at longer prediction horizons (336 and 720 steps), where the frequency branch helps maintain periodic pattern consistency. At horizon 720 on ETTh1, FTMamba achieves a 4.3\% MSE improvement over PatchTST, suggesting that frequency-domain modeling is particularly beneficial for long-term predictions.

\textbf{Dataset Characteristics.} FTMamba performs best on the ETTh1 and ETTh2 datasets, which contain electricity transformer temperature data with clear periodic patterns. The model shows relatively smaller improvements on ETTm1, where PatchTST performs competitively. This may be because ETTm1's 15-minute resolution introduces more fine-grained variations that are better captured by PatchTST's patch-based attention mechanism.

\subsection{Ablation Study}

To validate the contribution of each component, we conduct ablation experiments on ETTh1 with prediction horizon 96. The results are shown in Table~\ref{tab:ablation}.

\begin{table}[H]
\caption{Ablation study on ETTh1 (pred\_len=96).\label{tab:ablation}}
\begin{tabularx}{\textwidth}{CCC}
\toprule
\textbf{Configuration} & \textbf{MSE} & \textbf{MAE} \\
\midrule
FTMamba (Full)                         & \textbf{0.3826} & \textbf{0.4042} \\
w/o Frequency Branch (Mamba only)      & 0.3859   & 0.4114   \\
w/o Gated Fusion (simple addition)     & 0.3979   & 0.4149   \\
\bottomrule
\end{tabularx}
\end{table}

The results demonstrate that:
\begin{enumerate}
\item The frequency branch contributes to performance: removing it increases MSE by 0.9\% (0.3826 $\rightarrow$ 0.3859).
\item The gated fusion mechanism is critical: replacing it with simple addition increases MSE by 4.0\% (0.3826 $\rightarrow$ 0.3979), showing that adaptive fusion of temporal and frequency features is significantly better than naive combination.
\item The full FTMamba model achieves the best performance, validating the effectiveness of each component in the proposed architecture.
\end{enumerate}

%%%%%%%%%%%%%%%%%%%%%%%%%%%%%%%%%%%%%%%%%%
\section{Discussion}

The experimental results demonstrate that FTMamba effectively combines temporal and frequency-domain information for long-term time series forecasting. The consistent improvements across multiple datasets and prediction horizons validate the effectiveness of the proposed dual-branch architecture.

The gated fusion mechanism plays a crucial role in the model's performance. By learning to dynamically balance between temporal and frequency features, the model can adapt to different input characteristics and forecasting horizons. The 4.0\% MSE improvement over simple addition (Table~\ref{tab:ablation}) highlights the importance of adaptive fusion over naive combination.

The linear computational complexity of the Mamba architecture makes FTMamba scalable to long sequences, addressing a key limitation of Transformer-based methods. While PatchTST achieves competitive results on some horizons, FTMamba's frequency-aware design provides complementary benefits, particularly for longer prediction horizons where periodic patterns need to be maintained.

%%%%%%%%%%%%%%%%%%%%%%%%%%%%%%%%%%%%%%%%%%
\section{Conclusions}

We have presented FTMamba, a novel dual-branch architecture for long-term time series forecasting that combines Mamba's efficient temporal modeling with learnable frequency-domain feature extraction. The proposed gated fusion mechanism adaptively combines temporal and frequency information, enabling the model to leverage complementary patterns from both domains. Extensive experiments on three benchmark datasets demonstrate that FTMamba achieves state-of-the-art performance, outperforming strong baselines including PatchTST, iTransformer, and TimesNet across most prediction horizons.

Our work highlights the importance of frequency-domain modeling for time series forecasting and demonstrates that combining temporal and frequency information through learnable mechanisms can significantly improve prediction accuracy. The linear computational complexity of the Mamba architecture makes FTMamba scalable to long sequences, addressing a key limitation of Transformer-based methods.

Future work includes: (1) extending FTMamba to other time series tasks such as imputation and anomaly detection; (2) exploring more sophisticated frequency decomposition methods (e.g., wavelet transforms) for multi-scale analysis; and (3) investigating the model's interpretability by analyzing the learned frequency filters and gating patterns.

%%%%%%%%%%%%%%%%%%%%%%%%%%%%%%%%%%%%%%%%%%
\vspace{6pt}

%%%%%%%%%%%%%%%%%%%%%%%%%%%%%%%%%%%%%%%%%%
\authorcontributions{Conceptualization, D.P. and H.Y.; methodology, H.Y.; software, H.Y.; validation, D.P., H.Y. and H.J.; formal analysis, H.Y.; investigation, H.Y.; resources, D.P.; data curation, H.Y.; writing---original draft preparation, H.Y.; writing---review and editing, D.P.; visualization, H.Y.; supervision, D.P.; project administration, D.P.; funding acquisition, D.P. All authors have read and agreed to the published version of the manuscript.}

\funding{This research received no external funding.}

\institutionalreview{Not applicable.}

\informedconsent{Not applicable.}

\dataavailability{The data used in this study are publicly available from the Time-Series-Library repository at https://github.com/thuml/Time-Series-Library.}

\acknowledgments{The authors would like to thank the anonymous reviewers for their valuable comments and suggestions.}

\conflictsofinterest{The authors declare no conflicts of interest.}

%%%%%%%%%%%%%%%%%%%%%%%%%%%%%%%%%%%%%%%%%%
\begin{adjustwidth}{-\extralength}{0cm}

\reftitle{References}

\begin{thebibliography}{999}

\bibitem{ref1}
Rasp, S.; Thuerey, N. Data-driven medium-range weather prediction with a Resnet pretrained on climate simulations. \textit{J. Adv. Model. Earth Syst.} \textbf{2021}, \textit{13}, e2020MS002405.

\bibitem{ref2}
Hong, T.; Pinson, P.; Fan, S.; et al. Probabilistic energy forecasting: Global Energy Forecasting Competition 2014 and beyond. \textit{Int. J. Forecast.} \textbf{2016}, \textit{32}, 896--913.

\bibitem{ref3}
Li, Y.; Yu, R.; Shahabi, C.; Liu, Y. Diffusion convolutional recurrent neural network: Data-driven traffic forecasting. In Proceedings of the International Conference on Learning Representations (ICLR), Vancouver, BC, Canada, 30 April--3 May 2018.

\bibitem{ref4}
Zhang, Z.; Zohren, S.; Roberts, S. DeepLOB: Deep convolutional neural networks for limit order books. \textit{IEEE Trans. Signal Process.} \textbf{2020}, \textit{68}, 3515--3530.

\bibitem{ref5}
Vaswani, A.; Shazeer, N.; Parmar, N.; et al. Attention is all you need. In Proceedings of the Advances in Neural Information Processing Systems (NeurIPS), Long Beach, CA, USA, 4--9 December 2017; pp. 5998--6008.

\bibitem{ref6}
Wu, H.; Xu, J.; Wang, J.; Long, M. Autoformer: Decomposition transformers with auto-correlation for long-term series forecasting. In Proceedings of the Advances in Neural Information Processing Systems (NeurIPS), Virtual, 6--14 December 2021; pp. 22419--22430.

\bibitem{ref7}
Zhou, T.; Ma, Z.; Wen, Q.; et al. FEDformer: Frequency enhanced decomposed transformer for long-term series forecasting. In Proceedings of the International Conference on Machine Learning (ICML), Baltimore, MD, USA, 17--23 July 2022; pp. 27268--27286.

\bibitem{ref8}
Nie, Y.; Nguyen, N.H.; Sinthong, P.; Kalagnanam, J. A time series is worth 64 words: Long-term forecasting with transformers. In Proceedings of the International Conference on Learning Representations (ICLR), Kigali, Rwanda, 1--5 May 2023.

\bibitem{ref9}
Liu, Y.; Hu, T.; Zhang, H.; et al. iTransformer: Inverted transformers are effective for time series forecasting. In Proceedings of the International Conference on Learning Representations (ICLR), Vienna, Austria, 7--11 May 2024.

\bibitem{ref10}
Gu, A.; Dao, T. Mamba: Linear-time sequence modeling with selective state spaces. \textit{arXiv} \textbf{2023}, arXiv:2312.00752.

\bibitem{ref11}
Wang, S.; Zhang, H.; Li, L. S-Mamba: Mamba as a multivariate time series forecaster. \textit{arXiv} \textbf{2024}, arXiv:2403.11144.

\bibitem{ref12}
Chen, Y.; Li, Z.; Wang, S. TimeMachine: A time series is worth 4 mamba blocks. \textit{arXiv} \textbf{2024}, arXiv:2403.03820.

\bibitem{ref13}
Zeng, A.; Chen, M.; Zhang, L.; Xu, Q. Are transformers effective for time series forecasting? In Proceedings of the AAAI Conference on Artificial Intelligence, Washington, DC, USA, 7--14 February 2023; pp. 11121--11128.

\bibitem{ref14}
Wu, H.; Hu, T.; Liu, Y.; et al. TimesNet: Temporal 2D-variation modeling for general time series analysis. In Proceedings of the International Conference on Learning Representations (ICLR), Kigali, Rwanda, 1--5 May 2023.

\bibitem{ref15}
Gu, A.; Goel, K.; Ré, C. Efficiently modeling long sequences with structured state spaces. In Proceedings of the International Conference on Learning Representations (ICLR), Virtual, 25--29 April 2022.

\bibitem{ref16}
Yi, K.; Zhang, Q.; Fan, W.; et al. Frequency-domain MLPs are more effective learners in time series forecasting. In Proceedings of the Advances in Neural Information Processing Systems (NeurIPS), New Orleans, LA, USA, 10--16 December 2024.

\bibitem{ref17}
Zhou, H.; Zhang, S.; Peng, J.; et al. Informer: Beyond efficient transformer for long sequence time-series forecasting. In Proceedings of the AAAI Conference on Artificial Intelligence, Virtual, 2--9 February 2021; pp. 11106--11115.

\bibitem{ref18}
Wu, H.; Hu, T.; Liu, Y.; et al. Time-Series-Library: A comprehensive library for time series forecasting. Available online: https://github.com/thuml/Time-Series-Library (accessed on 1 May 2026).

\bibitem{ref19}
Kingma, D.P.; Ba, J. Adam: A method for stochastic optimization. In Proceedings of the International Conference on Learning Representations (ICLR), San Diego, CA, USA, 7--9 May 2015.

\end{thebibliography}

\end{adjustwidth}

\end{document}
